\chapter{Constructing and comparing competing interpretations}
\label{ch:search}

\section{The object of search is a legal choice}

Suppose three models explain paragraph 10 in different words. The first says that
product-linked gifts are prohibited unless they are fee discounts. The second
calls the same arrangement a marketing incentive and repeats the exception. The
third turns the sentence into an identical decision table. We have three
presentations but possibly only one interpretation. Their agreement tells us
little about the alternative that a mixed offer should be assessed component by
component, or about the evidence required to classify a cashback payment. An
ensemble that rewards the number of agreeing explanations may therefore become
more confident without becoming more informative.

A useful hypothesis changes a consequential choice. It might change the unit of
assessment from a whole offer to each benefit, distinguish reimbursement of a
paid charge from an unrelated reward, or require an incorporated definition
before applying a term. The hypothesis must also explain why that choice is
supported. Merely deleting the words about a product type creates a different
program, but an unsupported deletion is a test mutation, not an equally
respectable interpretation. The system needs both kinds of variation and must
label their roles accurately.

We will use a candidate record with five substantive parts: a proposition about
the source; the assumptions under which it holds; a formal rule that expresses
its consequences; supporting and opposing reasons; and a description of cases
that distinguish it from another candidate. The candidate also carries source
versions and an immutable identifier. This makes a claim such as ``this cashback
is a discount'' answerable: a reviewer can inspect the proposed definition,
identify the facts it requires, and observe the transactions on which a competing
definition would disagree.

The search objective follows from that representation. We seek a useful set of
supported alternatives, together with reasons for rejecting other proposals and
questions that remain unanswered. Finding a single high-scoring sentence is a
different objective. The former fits the bank's need to know where judgment is
required; the latter can conceal the judgment in an optimization procedure.

\section{Legal theories are more than predicted outcomes}

Theories of legal reasoning provide a richer starting point than generic text
search. \citet[sections 3--5]{theories2003} represent a theory through cases,
legally relevant factors, rules, preferences between rules, and values that
explain the preferences. A factor is a legally significant characteristic, not
simply a word that occurs in a judgment. A rule relates such characteristics to a
conclusion. A preference explains why one applicable rule should prevail over
another when both cannot determine the outcome together.

To understand the distinction, consider an invented training example about a
reimbursement policy. One decision allows reimbursement because the expense was
necessary for an authorized task. Another refuses it because the claimant had
already been reimbursed. A theory containing only the first rule predicts payment
whenever necessity is established. A theory containing both rules, but no
priority, encounters a conflict in the double-payment case. A theory in which the
bar on double payment defeats the necessity rule explains both outcomes. This is
a logical illustration, not a statement about an actual court or SFC policy.

The theory's values might explain a preference by reference to preventing double
recovery. That explanation is itself contestable: it must be connected to
applicable authority rather than invented as the bank's preferred commercial
outcome. The theoretical machinery helps represent such an explanation; it does
not decide which values an institution is legally entitled to prioritize.

For the circular, the analogous factors include a product link, a type-of-product
link, the character of the benefit, and evidence of a charge reduction. A theory
might treat an incentive's relationship to the subscription charge as decisive;
another might require that the charge itself is reduced in the contractual fee
schedule. Their predictions may coincide for a direct fee waiver and diverge for
cash paid after settlement. The divergence identifies an interpretive question
that a bare outcome label would obscure.

The source paper does not claim to recover complete judicial reasoning from raw
opinions. Section 5.4 explicitly discusses limits of the representation. Our
implementation must preserve that limit: converting a judgment to factors is
another interpretive task, with evidence spans, assumptions and review. A model
that extracts a convenient factor list and then proves a conclusion from it has
not bypassed the original source problem.

\section{What case-based argument contributes}

The history surveyed by \citet[sections 2--3]{hypolegacy2017} is particularly
relevant to the user's concern about a single reading. HYPO organizes exchanges
in which one side cites a favorable case, an opponent distinguishes it or presents
a counterexample, and the first side replies. This exchange encourages a search
for facts that make an apparent analogy fail. The purpose is not to collect the
largest number of vaguely similar documents.

Imagine that an earlier, applicable authority addresses a benefit described as a
fee rebate. A candidate relying on it must state the shared features: perhaps a
charge was first incurred, the later payment could not exceed that charge, and
both arose from the same transaction. An opposing candidate may distinguish an
offer in which the payment exceeds any charge or can be earned without incurring
one. These are questions about material similarities. A search engine's textual
similarity score does not answer them.

HYPO's claims lattice orders relationships between cases using legally relevant
dimensions. A lattice is a partial ordering: some comparisons remain
incomparable. Two precedents can each support a different aspect of the present
case, without either being globally ``more similar.'' Preserving this structure
is useful for an interpretation workbench. A numerical embedding can nominate
cases for inspection, while the argument record explains which feature makes a
case relevant and which feature might defeat the analogy.

CABARET, as described in the retrospective, combines statutory rules with
case-based reasoning about difficult terms. BankXX searches a network of pieces
of legal argument. These suggest a division of labor for the circular: apply
explicit logical structure directly, and investigate open-textured classifications
through sources and arguments. They do not justify allowing analogy to override
an unambiguous express restriction without an applicable legal reason. The
original CABARET and BankXX implementations have not been audited for this
monograph; the historical claims here come from the inspected retrospective.

A production case record therefore needs more than a citation and a summary.
Record the issuing body, jurisdiction, date, procedural context, proposition
relied on, relevant factual features, passage that supports the proposition, and
known treatment by later authority. Distinguish a court's holding from an
advocate's submission or an incidental observation. Those classifications must
remain reviewable. An invented or misclassified precedent is more dangerous when
attached to a polished formal argument than when presented as an uncertain search
result.

This book supplies no actual case-law ruling on the cashback example. Its role is
to specify how such authority would enter the system if found. The immediate
23EC46 exercise remains tied to the circular and its unresolved incorporated
materials. A reviewer should never have to infer that a hypothetical analogy is
a reported judicial decision.

\section{AGATHA: explicit construction of rival theories}

AGATHA supplies a direct precedent for searching over competing legal theories.
\citet[sections 2--4]{agatha2005} define moves that extend or challenge a developing
theory: analogy, distinction, counterexample and changes to preferences. Rather
than asking a model to ``think harder,'' the system asks what legal move would
change the current account and why. This makes the search history inspectable.

In our running example, an analogy constructor could propose that two payments
belong to the same classification because both reverse a charge. A distinction
constructor could split them by whether a charge was actually paid. An exception
constructor could attach the fee-discount qualification at a particular point in
the rule. A counterexample constructor could search for a scenario that the
candidate treats as permitted despite an express product-type link. Each move
should identify its parent, the changed proposition and the source support.

AGATHA evaluates theories using explanatory power, simplicity and completeness.
Its search is a modified A* procedure, and the paper explains why the modification
can miss the best theory. We should not import the word A* as an optimality
certificate. The evaluated cases already have factor descriptions. The reported
exercise is therefore substantially easier than discovering the factors,
resolving cross-references and determining current authority in raw banking
materials.

The paper's adversarial lookahead is also instructive. A theory can look strong
because a limited opponent fails to find the appropriate counterargument.
Optimizing for a victory against that opponent can favor a weak justification.
For a bank, the objective must be resistance to informed challenge, not persuasion
of a convenient critic. The system should preserve the strongest located
objection even when an advocate agent prefers to stop discussing it.

We borrow the explicit construction operators and the record of competing
reasons. We do not adopt the paper's heuristic weights as calibrated legal
confidence, nor its advocacy objective as the bank's release criterion. A
successful local adaptation must be evaluated on whether it finds material
alternatives and omissions within a declared source perimeter. That evaluation
has not yet been performed.

\section{Arguments, attacks and unresolved conclusions}

An argument is a chain from premises to a conclusion. Some steps are strict:
under the stipulated meaning of conjunction, two true components make the
conjunction true. Other steps are defeasible: evidence that a payment reverses a
charge may support a classification while remaining vulnerable to a contrary
source or missing condition. ASPIC+ separates these kinds of inference and
formalizes ways to challenge them \citep[sections 2--3]{aspic2010}.

There are three useful attacks. An attack on a premise disputes a starting fact,
such as whether the charge was paid. A rebuttal supports a contrary conclusion,
such as classifying the payment as an additional benefit. An undercut attacks the
inference itself: even if the payment and charge have equal amounts, equality of
amount alone may not establish that one legally discounts the other. These
attacks require different remedies. Retrieving a payment record may settle the
first; retrieving a definition may address the third. Repeating the same
classification prompt need not address either.

Let an argument graph contain arguments $a$ and $b$ that defeat each other, with
no stronger argument resolving the dispute. Under the usual preferred-extension
semantics, the graph has two alternative acceptable sets, $\{a\}$ and $\{b\}$.
Neither argument appears in both sets. Under grounded semantics, neither is
accepted in this simple graph. The result expresses unresolved competition; it
does not say both conclusions are legally true.

An extension is a set of arguments satisfying a specified acceptability
condition. A conclusion is credulously accepted when at least one extension
supports it, and skeptically accepted when every relevant extension supports it.
A conclusion may be supported by different arguments in different extensions;
requiring one identical argument to appear everywhere would be a stronger and
different condition. The selected semantics and treatment of premises must be
stored alongside the result.

There is an implementation trap at this point. If a semantics yields no
extensions, code that applies a universal predicate to an empty list may report
that every proposition passes. The inspected Carneades Dung implementation
illustrates why callers must handle the existence of extensions explicitly. The
workbench will require a nonempty computed collection before reporting skeptical
support. If no extension exists, or enumeration is incomplete, the result is a
separate unresolved status. It cannot become a positive conclusion through
vacuous truth.

Argumentation computes consequences of the supplied arguments, attacks and
preferences. It cannot establish that all relevant arguments were supplied or
that the legal preferences are justified. Moreover, ASPIC+ rationality results
depend on conditions on rules and knowledge bases; the author's correction to
definition 6.8 must accompany any implementation relying on those results. A
free-form graph of model-generated sentences does not automatically inherit the
paper's theorems.

\section{Burdens, assumptions and evidence in Carneades}

Carneades distinguishes ordinary premises, assumptions and exceptions
\citep[sections 3--6]{carneades2007}. An ordinary premise needs support. An
assumption may stand in a particular dialogue until challenged. An exception can
defeat an otherwise supported argument. These distinctions explain why absence
of evidence has different effects in different reasoning arrangements.

For example, the workbench might require affirmative evidence that a transaction
falls within the bank's activity profile. It must not silently assume that fact
because most transactions in a training set did. Conversely, a procedural review
may temporarily accept an explicitly declared factual assumption to compare the
consequences of two legal readings. That temporary acceptance is a property of
the review exercise, not a permission to use the assumed fact in production.

A burden of proof specifies what support is needed for a proposition in a given
context. The 2007 paper includes illustrative proof standards and explicitly
notes that they have not been validated as legal proof standards. It would be
wrong to rename one of its numerical or dialectical conditions ``SFC approval''
and install it as policy. The bank needs its own reviewed rules about which
questions require legal adjudication, what factual evidence is sufficient, and
which roles can accept residual operational risk.

The original model is acyclic. Newer Carneades software has a different
computational model and cannot be treated as an executable copy of every result
in the paper. Our source inspection of the newer implementation found useful
argument utilities but also potentially expensive extension enumeration. A
prototype should first support small, bounded graphs with explicit semantics and
known reference examples. Introducing a general argumentation engine before its
resource and semantic behavior are understood would create another opaque
component in the assurance chain.

For the first release, even a simple structured argument register is valuable.
Require each candidate to name its supporting premises, distinguish assumptions
from evidence, identify opponents and specify what would defeat it. This can be
implemented before automatic extension computation. The eventual engine then
formalizes an already visible review practice rather than replacing an undefined
practice with a numerical score.

\section{Breadth before expensive investigation}

The initial search should cover named dimensions: actor and activity scope,
subject identity, definitions, logical attachment of exceptions, modality,
quantification, time and incorporated context. These dimensions are prompts for
possible differences, not a Cartesian product of equally plausible answers. If a
choice has no source support, retain it as unsupported or as a negative control;
do not inflate the shortlist with implausible alternatives.

Consider three dimensions with two alternatives each. Exhaustive combination
would generate eight candidates, but some combinations may be inconsistent. A
whole-offer interpretation combined with a component-only factual schema may not
be executable without a transformation. A candidate that uses a definition from
one historical version and an exception from an incompatible later version may
not describe any applicable regime. Rejecting these combinations requires an
explicit reason, because an encoding defect could otherwise be mistaken for a
refutation of the underlying reading.

A practical initial frontier reserves room for each material family before
ranking refinements within a family. The reservation is a search policy, not a
claim that all families have equal legal merit. If the budget permits only six
candidates and eight supported families are identified, the report must say that
two families remain unexplored. It should not display a six-candidate shortlist
as exhaustive. Increasing a beam width can improve exploration, but no finite
width establishes semantic completeness for unrestricted legal language.

The search tree records how candidates were created. The argument graph records
shared dependencies. These are different structures. Two branches may rely on
the same disputed definition, and a later clarification should affect both.
Representing all dependencies only as parent-child edges would duplicate facts
and make invalidation unreliable. A directed acyclic graph of immutable candidate
revisions can coexist with a possibly cyclic argument graph, provided each
structure has its own semantics and resource limits.

Deduplication also needs layers. Exact duplicate bytes can be merged by hash.
Canonical syntax can identify expressions that differ only in irrelevant
formatting. Formal equivalence can establish equal decisions over a declared
domain. None of these establishes that source assumptions are identical. Two
candidates with equivalent Boolean bodies but different definitions of ``gift''
remain separate interpretive commitments. Their shared body may be stored once,
while their source and factual bindings remain distinct.

\section{Tree search and the meaning of a reward}

Tree of Thoughts separates candidate generation, evaluation and search
\citep[section 3]{treeofthoughts2023}. Its inspected breadth-first implementation
retains a beam selected by model-based values or votes. That is bounded search,
not exhaustive breadth-first enumeration. String deduplication in the code does
not ensure that the retained thoughts express different legal hypotheses.

Language Agent Tree Search combines tree search, model judgments, reflection and
environment feedback \citep[sections 3.2--4 and appendix A]{lats2024}. A commonly
used selection expression illustrates the tradeoff. For a parent visited $N$
times and child $i$ visited $n_i>0$ times, let $\bar r_i$ be its observed mean
reward. A UCT-style selector evaluates
\begin{equation}
 U_i=\bar r_i+c\sqrt{\frac{\log N}{n_i}},\qquad c>0.
 \label{eq:uct}
\end{equation}
The first term favors a child that has produced useful outcomes. The second
encourages investigation of less-visited children. A separate rule handles
unvisited children; setting their denominator to zero is not an implementation.
The constant $c$, reward definition and visit accounting are material design
choices and must be recorded.

For a game, the reward can be a terminal win. For a software task, a reward may
include tests. For interpretation, neither model agreement nor passing a
candidate's own tests is an external truth signal. A useful action reward could
instead record a verified dependency retrieved, a replayable disagreement
scenario found, or a previously open question resolved by authenticated evidence.
Such a reward measures progress in an investigation. It must not be displayed as
the probability that the resulting legal reading is correct.

The distinction is especially important when reflection uses the model's own
explanation of a failure. A model can learn to satisfy a flawed test or to agree
with its earlier answer. The resolution controller must keep the authoritative
question and evidence fixed while allowing the candidate to change. The oracle
cannot be rewritten merely to improve the candidate's reward.

For the first prototype, a deterministic scheduler with family coverage and
explicit action priorities is easier to validate than a learned MCTS policy.
MCTS can later be an optional scheduler under the same action, budget and release
constraints. A comparison would ask whether it finds more independently judged
material issues at equal cost. It would not ask whether it accumulates a larger
self-assigned score. Chapter~\ref{ch:evaluation} specifies the needed experiment.

\section{Choose questions that separate the surviving readings}

Suppose four supported candidates remain. Three agree on a direct fee waiver,
but they differ on a payment larger than the fee, a payment made without any fee,
and a mixed voucher-and-waiver offer. Rechecking the direct waiver provides
little discrimination. The system should identify the question whose answer is
most likely to separate the remaining positions.

Generalized binary search formalizes a related problem for a finite hypothesis
class and an answer oracle \citep[Figure 1 and sections I--II]{gbs2011}. Its
identification guarantees have assumptions about the hypothesis space, queries
and oracle. Legal interpretation may violate all three: a defensible reading may
not yet be generated, a question may not admit a definitive answer, and the
reviewer can be uncertain. We borrow question selection without claiming those
guarantees.

A transparent heuristic uses pairwise disagreement. Let $H$ be the retained
candidates, $w_h\geq0$ descriptive priority weights, and $o_h(x)$ the candidate's
predicted outcome for a proposed factual scenario $x$. Define
\begin{equation}
 D(x)=\sum_{\{h,j\}\subseteq H}w_h w_j\,
       \mathbf{1}\{o_h(x)\ne o_j(x)\}.
 \label{eq:disagreement-question}
\end{equation}
Each unordered pair contributes when the scenario makes its predictions differ.
With equal weights, this simply counts separated pairs. Dividing by the expected
cost of answering the question can help scheduling, but the cost estimate is
itself uncertain. Neither version converts the weights into posterior
probabilities.

For two candidates, the useful scenario is any feasible one on which they differ.
An SMT solver can sometimes construct such a scenario, subject to the domain
conditions explained in Chapter~\ref{ch:proof}. For classification disagreements,
the solver may expose a variable assignment but cannot manufacture evidence that
makes it legally meaningful. A witness with a ``discount'' flag set to true is not
an answer to the question of what establishes that flag.

The review question should therefore contain the facts, the rival conclusions,
the source passage responsible for the disagreement, and the additional evidence
or authority sought. ``Which answer is right?'' is usually too vague. ``Does the
incorporated definition require a reduction in the contractual charge, or can a
later reimbursement qualify when it is capped at the paid charge?'' identifies
the missing premise. If the available authority does not resolve it, that is a
substantive result that the report should preserve.

\section{Uncertainty about meaning is not a vote count}

Self-consistency samples multiple reasoning paths and aggregates their answers
\citep[section 2]{selfconsistency2023}. Its evaluations show benefits on the
studied reasoning benchmarks. The method deliberately distinguishes diversity of
paths from a single greedy trajectory, but the samples still share a model and
prompting context. Forty sampled answers are not forty independent legal experts.
The paper does not establish omission detection or calibrated legal confidence
for SFC circulars.

Semantic uncertainty goes further by grouping answers by meaning before
estimating uncertainty \citep[sections 3--4]{semantic2023}. This is helpful because
ten paraphrases of one answer should not look like ten separate answers. In the
idealized formulation, strings expressing the same meaning belong to a class
$c$, and the probability of the class is the sum of the probabilities of its
strings. The class entropy is then
\begin{equation}
 H(C)=-\sum_c p(c)\log p(c).
 \label{eq:semantic-entropy}
\end{equation}
A concentration on one class produces lower entropy than an even distribution
across several classes. This measures uncertainty of the model's answer
distribution under the selected grouping. It is not a proof that the concentrated
answer is correct.

The paper approximates semantic grouping with natural-language inference. Its
appendix discusses nontransitivity and sensitivity to clustering order. Our
inspection of the official source identified a further implementation issue: the
examined clustering condition tests the absence of label zero in either direction,
rather than explicitly requiring entailment in both directions. The retrieved model configuration maps label zero to contradiction, one to
neutral and two to entailment: the inspected code therefore also permits neutral
relations. The model configuration was retrieved separately from the pinned code;
its identity must be pinned as well before reuse. Even a
perfect implementation of the stated mutual-entailment test would remain an
empirical language-model judgment, not formal equivalence of bank controls.

The workbench can use semantic similarity to propose clusters for review. It
must not delete a dissenting candidate merely because an NLI model assigns a
high score. Formal equivalence can justify merging execution paths only within
its checked domain, while source assumptions and classification definitions stay
visible. Chapter~\ref{ch:evaluation} develops calibration and set-valued reporting
without assuming that a model's internal distribution is a distribution over
legal truth.

\section{What can be reused from the adjacent projects}

DynareMCP contains useful patterns for explicit hypotheses, diagnostics,
falsifiers and limits of claims. In the inspected legacy audit writer,
\srcpath{write_hypothesis_node} records a hypothesis, source support, code support,
a diagnostic, a falsifier and a supplied rank score. These fields map naturally
to interpretation candidates. The supplied score is not automatically calibrated;
its meaning depends on the caller.

The inspected legacy scheduler chooses among eligible frontier nodes by a stored
priority and identifier. It is not a general UCT engine hidden behind the word
``tree.'' Reusing its bookkeeping ideas is sensible; importing it and claiming
that a proven Monte Carlo search now verifies legal meaning would be wrong. The
new controller needs its own issue identity, source lineage, legal review states
and resource accounting.

MathDevMCP is useful once a proposition is precise enough to prove. A formal
obligation might concern preservation of a Boolean expression, absence of an
unbounded loop in a finite controller model, or equivalence of two encodings over
a declared domain. ResearchAssistant supports retrieval, parsing and retention
of the literature and regulatory evidence. Neither project supplies a privileged
interpretation of English. They strengthen different links in the chain.

The design that follows therefore combines several structures: a source inventory
to detect missing material, a candidate set to expose choices, an argument graph
to record reasons, a bounded search to investigate questions, and formal checks
to compare executable consequences. The diversity is useful because each
structure can reveal mistakes that the others overlook. Its value will depend on
how discrepancies move between them, which is the central subject of the next
chapter.
